\documentclass{article}
\setlength\parindent{0pt}
\setlength\parskip{6pt}

\usepackage[backend=biber]{biblatex}
\usepackage{xcolor}

\title{``Some ideas"}
\author{}

\addbibresource{biblio.bib}


\begin{document}
\maketitle


\section{Parkinson's}
Parkinson's desease is the second most diffuse neurodegenerative disorder
originating from the progressive degeneration of key dopaminergic nuclei in the
brain, and leading to a progressive impairment of the patients motor and
cognitive abilities \cite{Pringsheim2014ThePO}.

Although a considerable amount of effort and resources has been put into the
study of this condition and of the possible ways to undermine its developement,
a properly disease-modifying treatment has not yet been reached [some ref], and
most of the therapeutic options currently available can only act by limiting
the symptoms or slowing down the desease progression \cite{Marras2020TherapyOP}.

In order to maximize the therapeutic effect of these therapeutic options, as well
as study their potential early-stage disease-modifying potential, earyly diagnosis
is fundamental \cite{Pagn2012ImprovingOT}. 

For this purpose, numerous pre-motor signs have been  identified (e.g. 
idiopathic RBD, constipation) and represent classically known risk factors for 
Parkinson's developement [some book, this is old stuff]. However, identifying and
correctly interpreting these symptomatic constallations is difficult, and a clear,
quantitative, measure related to a clear desease sign ({\color{blue} qui il 
biomarker?}) is still lacking. 

Recent studies have made remarkable steps in this direction, leveraging in
particular, the non-invasive and yet highly-throuput data-generating device
that are the smartdevices. In particular, it has been shown that a number of
gait-related quantitative metrics can be derived from this type of data
\cite{Reimer2022Evaluating3H}, and used to quantify Parkinson's deasease
progression and onset \cite{Shahar2021GaitAU}. 

These high-volume datasets are, moreover, exceptionally suitable for the developement
of novel data-driven approaches and, in particular, machine-learning and deep-learning
guided analysis \cite{Mei2021MachineLF}. 

{--\color{red} Qui secondo me parte la parte che si vuole proporre per il grant--}

Multi-modal data, generated by the simulataneous acquisitions of a smartphone, {\color{red}
il vostro device che fa gli ECG e poi non so qualcos'altro} describes, in particular,
the temporal evolution of an underlying biological network, and this is the natural 
framework in which Geometric Deep Learning algorythms and Graph Neural Networks can be used.









\printbibliography

\end{document}
